% Section 6.1, from lectures/L06/appendix/a_driver_stack.md.

\section{The Driver Stack (the Contracts)}
\label{c6:sec:driver_stack}\label{c6:app:a}
This chapter crosses from VHDL to C++, and it starts with the \emph{shapes}: the three abstractions
the driver will rest on. None of them contains an algorithm yet; they are the register map, the
transport interface, and the driver interface. Designing them well is what makes the driver
(\chapref{7}) short and host-testable.

Two of the three are abstract interfaces, so each lives in its own namespace: the driver's public
API is \code{driver::uart::Interface}, and the transport seam it runs over is
\code{driver::transport::Interface}. Naming both \code{Interface}, in separate namespaces, is
deliberate: each subsystem exposes one seam by that name.

Everything here is written by you. The three contracts are declarations; the exercises then
implement the driver interface once, as \code{driver::uart::Stub}, so the chapter ends with a
concrete class rather than three headers. The provided host tests arrive in \chapref{8}, once the
\code{Uart} that implements the same interface exists.

The full stack, top to bottom:

\begin{booktable}{@{}>{\raggedright\arraybackslash}p{15.2em}>{\raggedright\arraybackslash}p{3.8em}L@{}}
  \thead{Layer} & \thead{Chapter} & \thead{Role} \\
  \midrule
  The application, \code{app::EchoNode} & \ref{c10:ch} & Uses the driver through its abstract
    interface. \\
  \code{driver::uart::Interface} & \ref{c6:ch} & The abstract driver API the application depends
    on. \\
  \code{driver::uart::Stub} & \ref{c6:ch} & A test double implementing that interface; the whole
    harness for \code{app::EchoNode} in \chapref{10}. \\
  \code{driver::uart::Uart} & \ref{c7:ch} & Implements the driver interface using the register
    protocol. \\
  \code{driver::transport::Interface} & \ref{c6:ch} & The abstract SPI byte seam the driver runs
    over. \\
  A scripted stub, or the real SPI & \ref{c8:ch}, \ref{c9:ch} & \code{driver::transport::Stub} for
    the host tests; the AVR SPI on hardware. \\
\end{booktable}

\noindent Every layer is C++. The register map both halves share lives in
\file{driver/uart/register_map.hpp}.

\subsection{The register map}
\label{c6:sec:register_map}
The C++ side needs the same register map the hardware has: the seven register indices and the
meaning of each \code{STATUS}, \code{CTRL} and \code{ERROR_FLAGS} bit, from Part~2 of the protocol
(\secref{proto:part2}). It comes from the same source as \file{uart_def.vhd} on the other side of
the wire, and where a name exists on both sides the value must be identical; you declare only the
bits this driver actually uses. You transcribe it once, as named constants, so that nothing above
ever writes a bare \code{3} where it means ``the \code{TX_DATA} register''. A value here that
disagrees with the hardware is a bug you would only find on the bench, so it is worth copying
carefully.

\subsection{The transport seam}
\label{c6:sec:transport_seam}
The driver will never talk to SPI directly. It talks to an abstract seam,
\code{driver::transport::Interface}, that does one thing: exchange raw bytes over one SPI
transaction. A transaction is \emph{select} (pull \code{SS} low), some byte exchanges, then
\emph{deselect} (release \code{SS}), matching the framing in Part~3 of the protocol
(\secref{proto:part3}). Each byte exchange is duplex: the same call sends one byte and returns the
byte that arrived on \code{MISO} at the same time, because that is exactly how SPI works, a bit out
and a bit in on every clock.

Putting the seam at the \emph{byte} level, rather than exposing register reads and writes directly,
is the deliberate choice. It keeps the SPI-specific knowledge (five-byte transactions, the command
byte, the byte order) in one place, the driver, and lets the layer below be a dumb byte pipe that a
stub can imitate perfectly. The seam is the whole reason the driver is host-testable: the tests of
\chapref{8} inject a stub here, \chapref{9} injects the real AVR SPI, and nothing above the seam
changes.

\subsection{The driver interface}
\label{c6:sec:driver_interface}
\code{driver::uart::Interface} is the abstract API the application codes against, and it promises
five things: configure the peripheral, write a byte, read a byte, report the status and error
state, and clear errors. Making it abstract is what lets the application, and its tests, be written
against a promise rather than a concrete driver, and it is why the same application later runs
unchanged over the real hardware driver.

\code{write()} and \code{read()} are both \textbf{non-blocking}. \code{write()} pushes a byte only if
the transmitter has room, and reports whether it did; \code{read()} hands back a byte only if one
has arrived. That maps directly onto the FIFO-backed \code{STATUS} bits of \chapref{5} and makes
every operation deterministic to test. Blocking versions, which spin until the peripheral is ready,
are a thin convenience built on top, an exercise in \chapref{8} rather than part of the core.

\subsection{What's ahead}
The exercises in \secref{c6:sec:exercises} build the register map, the transport interface and the
driver interface, then implement that interface once as \code{driver::uart::Stub}. In \chapref{7}
you implement the \code{Uart} driver over these contracts, using the register protocol, and in
\chapref{8} you test it with a scripted \code{driver::transport::Stub}.
